\chapter{Integrale di Lebesgue}

Finalmente arriviamo a uno dei capitoli più fecondi e profondi della teoria della misura e dell’integrazione: l’Integrale di Lebesgue.\\
\\
In precedenza abbiamo costruito con cura lo spazio di probabilità $(\Omega, \mathcal{A}, \mathbb{P})$ e approfondito i concetti di misura, sigma-algebra e variabile aleatoria. Ora entriamo in campo con uno strumento che supera i limiti della tradizionale integrazione di Riemann: consente di trattare funzioni più “irregolari”, serie di funzioni, scambi di limite e integrale, e fornisce un linguaggio comune per analisi matematica, probabilità e teoria delle funzioni.\\
\\
In questo capitolo vedremo come definire l’integrale di una funzione misurabile rispetto a una misura — in particolare la misura di probabilità — e illustreremo perché la sua struttura (attraverso funzioni semplici, limiti monotoni, e la proprietà «quasi ovunque») risulti essenziale nella moderna teoria della probabilità e dell’analisi.

\medskip
\noindent Abbandonando la sola visione geometrica ``area sotto la curva'' tipica dell’integrazione di Riemann, l’approccio di Lebesgue ci mostra come dividere il codominio della funzione in “strati” orizzontali e misurare quanto si spende in ciascuno di essi. Questo cambio di prospettiva rende l’integrazione molto più flessibile, ci consente di definire il valore atteso $E[X]$ di variabili aleatorie in modo robusto, e di applicare teoremi fondamentali come il teorema della convergenza dominata o monotona.

\medskip
\noindent Con questo nuovo strumento potremo affrontare variabili che prima risultavano fuori dalla portata, funzioni definite su spazi astratti, e avanzare verso gli spazi $\mathcal{L}^p$ e le applicazioni in probabilità e analisi funzionale.
Prepariamoci dunque a entrare nel cuore dell'argomento: comprenderemo prima la definizione di integrale per funzioni semplici non-negative, fino ad estenderlo a funzioni generalizzate, e vedremo le sue proprietà essenziali.

\section{Definizione per Funzioni Elementari}
D'ora in avanti, per rimanere più generali possibili, consideriamo $(E, \mathcal{E}, \mu)$ spazio di misura, e supponiamo che $\mu$, sia una misura $\sigma$--finita\footnote{$\mathbb{P}$ è una misura $\sigma$--finita.} (Definizione \ref{def:sigmafinita}).

\begin{definizione}[Funzione Elementare]
    $f:E\to(0, +\infty)$ è una funzione elementare se:
    \[\exists n\geq1, \quad A_1, ..., A_n\subseteq\mathcal{E}\;\text{ t.c. }\;A_h\cap A_k = \varnothing\;\;\forall h\neq k, \quad \bigcup_{k=1}^n A_k=E, \quad a_1...a_n\geq 0\;\text{tale che:}\]
    \[\boxed{f(x) = \sum_{k=1}^n a_k \;\mathbb{I}_{A_k}(x)}\]
\end{definizione}
A prima vista può sembrare una definizione complicata, mentre in realtà è abbastanza semplice. Supponiamo per comodità che $n=6$, allora possiamo partizionare $\mathcal{E}$ in sei diversi sottoinsiemi (disgiunti). La nostra funzione in quegli insiemi prende dei valori, che sono $a_1, a_2, ...,a_6$.

\begin{figure}[H]
\centering

% --- Primo disegno (rettangolo con a_i) ---
\begin{minipage}{0.48\textwidth}
\centering
\begin{tikzpicture}[x=0.8cm, y=0.8cm, line cap=round, line join=round, font=\small]

% --- Rettangolo esterno (Omega) ---
\draw[very thick] (0,0) rectangle (8,3);
\node[right] at (8.1,1.5) {$\boldsymbol{\mathcal{E}}$};

% --- Curve interne (sottoinsiemi) ---
\draw[smooth, thick]
  (0.5,0).. controls (0.8,1.2) and (1.2,2.3)..(1.6,3);

\draw[smooth, thick]
  (2.2,0).. controls (2.4,1.0) and (2.6,2.3)..(2.9,3);

\draw[smooth, thick]
  (4.0,0).. controls (4.2,1.2) and (4.5,2.5)..(4.8,3);

\draw[smooth, thick]
  (2.2,0).. controls (3,2.3) and (3.2,1.0)..(5.5,1.7)
  .. controls (6.5,2.0) and (7.3,1.4)..(7.5,0);

% --- Etichette dei sottoinsiemi ---
\node[red] at (0.8,2.2) {$a_1$};
\node[red] at (1.9,2.3) {$a_2$};
\node[red] at (3.3,2.3) {$a_3$};
\node[red] at (5.9,2.4) {$a_4$};
\node[red] at (3.8,0.8) {$a_5$};
\node[red] at (6.8,0.7) {$a_6$};
\end{tikzpicture}
\end{minipage}
\hfill
% --- Secondo disegno (asse reale con A_i) ---
\begin{minipage}{0.48\textwidth}
\centering
\begin{tikzpicture}[x=0.8cm,y=0.8cm,>=stealth,font=\small]

% ---- parametri dei tagli ----
\def\a{-2}
\def\b{2}

% ---- assi ----
\draw[->,thick] (-5,0) -- (5,0) node[right] {$\mathbb{R}$};
\draw[->,thick] (0,-0.2) -- (0,3) node[above] {$f(x)$};

% tacche e 0
\foreach \x in {-4,-3,-2,-1,1,2,3,4}
  \draw (\x,0.08) -- (\x,-0.08);

% ---- A2: (a,b) ----
\node[orange] at ({(\a+\b)/2},-0.55) {$A_2$};
\node[orange] at (\a,0) {\large $($};
\node[orange] at (\b,0) {\large $)$};

% segmento rosso sopra A2 e label a2
\draw[red,very thick] (\a,1.2) -- (\b,1.2);
\node[red] at ({0.3},1.55) {$a_2$};

% ---- punti rossi di esempio su A1 e A3 ----
\fill[red] (\a-1,1.2) circle (1.6pt) node[left=2pt,red] {$a_1$};
\fill[red] (\b+1,1.2) circle (1.6pt) node[right=2pt,red] {$a_3$};

\draw[red,very thick] (-4,0) -- (-3.05,0);
\draw[red,very thick] (-2.95,0) -- (-2.05,0);
\draw[red,very thick] (2.05,0) -- (2.95,0);
\draw[red,very thick] (3.00,0) -- (4,0);

\draw[red,very thick, dashed] (-4.5,0) -- (-4,0);
\draw[red,very thick, dashed] (4,0) -- (4.5,0);

% ---- A3: [b,+inf) ----
\fill[orange] (3,0) circle (2pt);
\node[orange] at (3,-0.55) {$A_3$};

% ---- A1: (-inf, a] ----
\fill[orange] (-3,0) circle (2pt);
\node[orange] at (-3,-0.55) {$A_1$};

\end{tikzpicture}
\end{minipage}

\caption{Prima un esempio di con $E\neq\mathbb{R}$, poi con $E=\mathbb{R}$}
\end{figure}


Se semplifichiamo maggiormente le cose, e consideriamo il caso in cui $E=\mathbb{R}$, allora il grafico sulla destra esprime perfettamente il concetto di funzione semplice.

\begin{definizione}
\label{def:intleb}
    Per ogni $f\geq0$ elementare si definisce l'Integrale di Legesgue di $f$ su $E$ rispetto alla misura $\mu$ il seguente numero (oppure $+\infty$)\footnote{Potrebbe valere $+\infty$ siccome $\mu$ è una misura $\sigma$--finita, e non finita}:
    \[\boxed{\int_E f\;d\mu:= \sum_{k=1}^n a_k \; \mu(A_k)}\quad \geq 0\]
\end{definizione}

\begin{proposizione}
    L'integrale di Lebesgue è lineare e monotono, ossia:
    \begin{itemize}
        \item siano $\;a, b\geq0; \quad f, g\in \mathbb{E}^+\;$ (ossia lo spazio delle funzioni elementari positive):
        \[\Rightarrow\quad af+bg\in\mathbb{E}^+\quad \text{e}\quad \int_E(af+bg)\;d\mu = a\int_Ef\;d\mu + b\int_E g\;d\mu\]
        \item siano $f, g\in\mathbb{E}^+;\quad f\leq g$:
        \[\Rightarrow\quad \int_Ef\;d\mu\leq\int_Eg\;d\mu\]
    \end{itemize}
\end{proposizione}

\section{Generalizzazione della Definizione a Funzioni Generiche}
\noindent Ora che abbiamo definito l'integrale per funzioni elementari positive, vogliamo chiaramente cercare di estendere la definizione in maniera rigorosa a una funzione generale. 

\medskip
\noindent Iniziamo vedendo che possiamo scrivere una qualsiasi funzione positiva, come una successione di funzioni in $\mathbb{E}^+$, potendo così estendere facilmente la Definizione \ref{def:intleb} a questa nuova classi di funzioni.
\begin{proposizione}
\label{prop:6.4}
    Per ogni $f:E\to[0, +\infty], \;\exists\;$ una successione $f_n:E\to[0+\infty)$ tale che:
    \begin{enumerate}
        \item $f_n\in\mathbb{E}^+, \quad \forall n\geq1$
        \item $f_n\uparrow f\;$, ossia che:
        \(\qquad \quad \underset{n\to+\infty}{\lim}f_n(x) = f(x) \qquad \text{e che}\qquad f_{n+1} \geq f_n\;\;\forall x\in E\)
    \end{enumerate}
\end{proposizione}

\newpage
\begin{definizione}
\label{def:6.5}
    Per ogni $f:E\to [0, +\infty]$ l'Integrale di Legesgue di $f$ su $E$ rispetto alla misura $\mu$  è  dato da:
    \[\boxed{\int_E f\;d\mu = \lim_{n\to+\infty}\int_E f_n\;d\mu}\]
\end{definizione}

\bigskip
\noindent Facciamo alcune osservazioni fondamentali.\\
\\
\(\bigstar\) Siccome valgono entrambi i punti 1. e 2. della Proposizione \ref{prop:6.4}, allora avremo che: $\int_E f_{n+1}\;d\mu\geq\int_E f_n\;d\mu$.
\begin{align*}
    &\Rightarrow\quad \left(\int_E f_n\;d\mu\right)_{n\geq1}\; \text{ è  una successione monotona non decrescente}\\
    &\Rightarrow\quad \exists\;\lim_{n\to+\infty}\int_E f_n\;d\mu = \begin{cases}
    0\leq L<+\infty\\
    +\infty
    \end{cases}
\end{align*} 
Il limite non può essere $<0$, in quanto $f_n\in\mathbb{E}^+, \;$ quindi al limite può valere $0$, ed è chiaro che $\int_E 0\;d\mu = 0$\\
\\
\(\bigstar\) Si può dimostrare che il limite presente nella Definizione \ref{def:6.5} non dipende dalla particolare successione approssimante $f_n$ che si usa per approssimare $f$, ovvero che:
\[\text{Se: }\quad \exists \tilde{f_n}\uparrow f\qquad\Rightarrow\qquad \lim_{n\to+\infty}\int_E \tilde{f_n}\;d\mu = \lim_{n\to+\infty}\int_E f_n\;d\mu\]
Di conseguenza $\int_E f\;d\mu$ risulta essere ben definito.\\
\\
\(\bigstar\) Potremmo legittimamente chiederci quale possa essere una successione $f_n$ che riesca ad approssimare crescendo $f$. 

\medskip
\noindent Costruiamo la seguente successione di funzioni, per ogni $n\geq1\;\text{ e } \;k = 0, ... , n\cdot2^n-1$, definita come segue:
\[
f_n(x) =
\begin{cases}
\dfrac{k}{2^n}, \quad & \text{se } \; \dfrac{k}{2^n} \le f(x) < \dfrac{k+1}{2^n},\\[6pt]
n, & \text{se } \;f(x) \ge n
\end{cases}
\]
Notiamo già il principio di quella che è stata l’intuizione fondamentale di 
Lebesgue: nella costruzione di questa successione, l’attenzione non è posta sul \textit{dominio} della funzione, come avveniva nell’integrale di Riemann, ma sul suo \textit{codominio}.

\medskip
\noindent Nel metodo di Riemann, infatti, si suddivide l’intervallo di definizione della funzione (cioè il dominio) in piccoli sottointervalli, e si sommano le aree dei rettangoli aventi base in ciascun sottointervallo e altezza data dal valore della funzione in un punto di esso.

\medskip
\noindent L’idea di Lebesgue, invece, è radicalmente diversa: invece di “tagliare” l’asse delle ascisse, si suddivide l’asse delle ordinate, cioè l’insieme dei valori che la funzione assume. In questo modo si costruiscono, per ogni livello del codominio, gli insiemi di punti del dominio in cui la funzione assume valori compresi in un certo intervallo, e si considera la misura di questi insiemi.

\medskip
\noindent Questa inversione di prospettiva  è ciò che permette all’integrale di Lebesgue di trattare funzioni molto più generali, di gestire limiti di successioni di funzioni e di garantire risultati di convergenza (che vedremo) che l’integrale di Riemann non poteva assicurare.\\
\\
Per essere più chiari, fissiamo $n=1$. Allora avremo che:

\[
f_1(x) =
\begin{cases}
0 \;(=a_1)
    &\quad \text{se } 0\le f(x) < \tfrac{1}{2} \quad (= A_1),\\[6pt]
\tfrac{1}{2} \;(=a_2) 
    &\quad \text{se } \tfrac{1}{2} \le f(x) < 1 \quad (= A_2),\\[6pt]
1 \;(=a_3)
    &\quad \text{se } f(x) \ge 1 \quad (= A_3).
\end{cases}
\]

\begin{figure}[H]
\centering
\begin{tikzpicture}[x=1.5cm,y=2cm,>=stealth,font=\small]

% --- Assi
\draw[->,thick] (-4.7,0) -- (4.7,0) node[right] {$x$};
\draw[->,thick] (0,0) -- (0,1.8) node[above] {$f(x)$};

% --- Soglie orizzontali
\draw[dashed] (-4.6,1) -- (4.6,1) node[right] {$a_1 = 1$};
\draw[dashed] (-4.6,0.5) -- (4.6,0.5) node[right] {$a_2 = \tfrac12$};

% --- Funzione : f(x) = 0.75 + 0.5 sin(60 x) + 0.15 sin(2 x)
\draw[thick,smooth,domain=-4.55:4.55,samples=220]
  plot(\x,{0.75 + 0.5*sin(60*\x) + 0.15*sin(2*\x r)});

% --- Punti verticali ESATTI come richiesto
\foreach \x in {-3.9,-2.15,-0.3,0.3,2.15,3.9}{
  \draw[dotted] (\x,0) -- (\x,1.4);
}

% ==================== f_1(x) a gradini, determinato automaticamente ====================

% Intervalli: I1=[-4.6,-3.9], I2=[-3.9,-2.15], I3=[-2.15,-0.3],
%             I4=[-0.3,0.3], I5=[0.3,2.15], I6=[2.15,3.9], I7=[3.9,4.6]

% Macro per disegnare un gradino sul livello corretto:
% usa il valore di f a metà dell'intervallo per scegliere 0 / 0.5 / 1
\newcommand{\stepfromto}[2]{%
  \pgfmathsetmacro{\xa}{#1}
  \pgfmathsetmacro{\xb}{#2}
  \pgfmathsetmacro{\xm}{(\xa+\xb)/2}
  \pgfmathparse{0.75 + 0.5*sin(60*\xm) + 0.15*sin(2*\xm r)}\let\val\pgfmathresult
  \pgfmathsetmacro{\yl}{(\val>=1 ? 1 : (\val>=0.5 ? 0.5 : 0))}
  \draw[orange,very thick] (\xa,\yl) -- (\xb,\yl);
}

% Orizzontali (sette intervalli)
\stepfromto{-4.6}{-3.9}
\stepfromto{-3.9}{-2.15}
\stepfromto{-2.15}{-0.3}
\stepfromto{-0.3}{0.3}
\stepfromto{0.3}{2.15}
\stepfromto{2.15}{3.9}
\stepfromto{3.9}{4.6}


%etichette A1/A2/A3:
\node[below] at (-4.25,0) {$A_3$};
\node[below] at (-3.025,0) {$A_2$};
\node[below] at (-1.225,0) {$A_1$};
\node[below] at (0, 0) {$A_2$};
\node[below] at (4.25,0) {$A_3$};
\node[below] at (3.025,0) {$A_2$};
\node[below] at (1.225,0) {$A_1$};
\node[orange, above] at (2.7,1.2) {$f_1(x)$};

\end{tikzpicture}
\end{figure}

\noindent Rimarchiamo che il dominio viene suddiviso in funzione del codominio, non il contrario.\\
Andando avanti con il ragionamento, vediamo come si comporta la nostra successione con $n=2$. Avremo che:

\[
f_{2}(x)=
\begin{cases}
0,          & \text{se } 0 \le f(x) < \tfrac14,\\[3pt]
\tfrac14,   & \text{se } \tfrac14 \le f(x) < \tfrac12,\\[3pt]
\tfrac12,   & \text{se } \tfrac12 \le f(x) < \tfrac34,\\[3pt]
\tfrac34,   & \text{se } \tfrac34 \le f(x) < 1,\\[3pt]
1,          & \text{se } 1 \le f(x) < \tfrac54,\\[3pt]
\tfrac54,   & \text{se } \tfrac54 \le f(x) < \tfrac32,\\[3pt]
\tfrac32,   & \text{se } \tfrac32 \le f(x) < \tfrac74,\\[3pt]
\tfrac74,   & \text{se } \tfrac74 \le f(x) < 2,\\[3pt]
2,          & \text{se } f(x) \ge 2.
\end{cases}
\]
\noindent Il cui grafico risulta essere:


\begin{figure}[H]
\centering
\begin{tikzpicture}[x=1.5cm,y=2cm,>=stealth,font=\small]

% --- Assi
\draw[->,thick] (-4.7,0) -- (4.7,0) node[right] {$x$};
\draw[->,thick] (0,0) -- (0,1.8) node[above] {$f(x)$};

% --- Soglie orizzontali (k/4, k=1..8)
\foreach \k/\lab in {1/{\tfrac14},2/{\tfrac12},3/{\tfrac34},4/{1},5/{\tfrac54}}{
  \draw[dashed] (-4.6,\k/4) -- (4.6,\k/4);
}

\node[right] at (4.6,0.25) {$a_2=\tfrac14$};
\node[right] at (4.6,0.5) {$a_3=\tfrac12$};
\node[right] at (4.6,3/4) {$a_4=\tfrac34$};
\node[right] at (4.6,1) {$a_5=1$};
\node[right] at (4.6,5/4) {$a_6=\tfrac54$};


% --- f(x)
\draw[thick,smooth,domain=-4.55:4.55,samples=220]
  plot(\x,{0.75 + 0.5*sin(60*\x) + 0.15*sin(2*\x r)});

% --- f_2(x): quantizzazione a passi di 1/4, troncata tra 0 e 2
  \def\xmin{-4.55}
  \def\xmax{4.55}
  \def\samples{1600}
  \pgfmathsetmacro{\dx}{(\xmax-\xmin)/\samples}

  % --- Funzione base e quantizzazione a gradini
  % f_raw(x) = 0.75 + 0.5 sin(60 x) + 0.15 sin(2 x)    (2x in radianti)
  % f_2(x)   = min(2, 0.25 * floor(4 * max(0, f_raw(x))))
  \pgfmathsetmacro{\x}{\xmin}
  \pgfmathsetmacro{\yraw}{0.75 + 0.5*sin(60*\x) + 0.15*sin(2*\x r)}
  \pgfmathsetmacro{\y}{min(2, 0.25*floor(4*max(0,\yraw)))}

  % Memorizza inizio gradino
  \edef\xstart{\xmin}
  \edef\yconst{\y}

  % --- Scansione e disegno SOLO dei segmenti orizzontali
  \foreach \i in {1,...,\samples}{
    \pgfmathsetmacro{\xnext}{\xmin + \i*\dx}
    \pgfmathsetmacro{\yrawnext}{0.75 + 0.5*sin(60*\xnext) + 0.15*sin(2*\xnext r)}
    \pgfmathsetmacro{\ynext}{min(2, 0.25*floor(4*max(0,\yrawnext)))}
    % Se cambia il livello, chiudi il gradino corrente (solo orizzontale) e riparti
    \pgfmathparse{(\ynext==\yconst)?1:0}
    \ifnum\pgfmathresult=1
      % nulla, continua
    \else
      % disegna il segmento orizzontale del gradino appena concluso
      \draw[very thick,red] (\xstart,\yconst) -- ({\xnext-\dx},\yconst);
      % aggiorna inizio e livello del nuovo gradino
      \xdef\xstart{\xnext}
      \xdef\yconst{\ynext}
    \fi
  }

  % --- Disegna l'ultimo gradino fino a xmax
  \draw[very thick,red] (\xstart,\yconst) -- (\xmax,\yconst);

  % --- Etichetta
  \node[red,above] at (2.7,1.3) {$f_2(x)$};

% Intersezioni di f(x)=0.75 + 0.5*sin(60x) + 0.15*sin(2x r) con a_k = k/4
% (dominio x ∈ [-4.55, 4.55])

% a = 1/4
\draw[dotted] (-1.56673,0) -- (-1.56673,1.4);
\draw[dotted] (-0.74196,0) -- (-0.74196,1.4);

% a = 1/2
\draw[dotted] (-2.15255,0) -- (-2.15255,1.4);
\draw[dotted] (-0.31455,0) -- (-0.31455,1.4);
\draw[dotted] (3.88372,0) -- (3.88372,1.4);

% a = 3/4
\draw[dotted] (-2.83368,0) -- (-2.83368,1.4);
\draw[dotted] (0.00000,0) -- (0.00000,1.4);
\draw[dotted] (2.83368,0) -- (2.83368,1.4);

% a = 1
\draw[dotted] (-3.88372,0) -- (-3.88372,1.4);
\draw[dotted] (0.31455,0) -- (0.31455,1.4);
\draw[dotted] (2.15255,0) -- (2.15255,1.4);

% a = 5/4
\draw[dotted] (0.74196,0) -- (0.74196,1.4);
\draw[dotted] (1.56673,0) -- (1.56673,1.4);


\end{tikzpicture}
\end{figure}

\noindent Vediamo graficamente che più $n$ aumenta, più $f_n$ si avvicina ad $f$. Inoltre sia $f_1$ che $f_2$ risultano essere funzioni elementari. Questa non è una dimostrazione rigorosa che $f_n\uparrow f$ e che $f_n\in\mathbb{E}^+$, tuttavia ci mostra intuitivamente come funziona il processo ``avvicinamento'' da parte della successione verso la funzione.\\
\\
\(\bigstar\) Restano vere le proprietà di linearità e monotonia anche per l'integrale definito per funzioni qualsiasi positive, ossia che:
$\forall a, b >0, \;\forall f, g:E\to[0, +\infty], \;$ allora si ha che:
\begin{itemize}
    \item \(\int_E(af+bg)\;d\mu = a\int_Ef\;d\mu + b\int_E g\;d\mu\)
    \item \(0\leq f\leq g\quad \Rightarrow\quad 0\leq\int_Ef\;d\mu\leq\int_Eg\;d\mu\)
\end{itemize}

\noindent Come abbiamo detto in precedenza, per l'integrale di Lebesgue riusciamo a far valere dei criteri di convergenza che non valevano per l'integrale di Reimann:
\begin{teorema}[Beppo Levi o Convergenza Monotona]
    Sia $f_n\uparrow f, \; f_n\geq 0\;\forall n\geq1, \; f\geq0$, allora vale che:
    \[
    \int_E f_n\,d\mu \;\uparrow\; \int_E f\,d\mu.
    \]
\end{teorema}

Questo risultato è quindi caratteristico dell'integrale di Lebesgue e \emph{non vale} in generale per l'integrale di Riemann. 

\medskip
\noindent Infatti, l'integrale di Riemann non gode di buone proprietà di continuità rispetto alla convergenza puntuale: per scambiare limite e integrale è necessario che la convergenza sia \emph{uniforme}. 

\medskip
\noindent Intuitivamente si può pensare alla convergenza monotona, come una convergenza dove ogni ``grafico'' $f_n$ si avvicina a $f$ \textit{punto per punto}, ma in modo \emph{indipendente} per ciascun $x$; la convergenza può essere più lenta o più rapida a seconda del punto considerato. Invece per la convergenza monotona si ha che l'intero grafico di $f_n$ si ``schiaccia'' su quello di $f$ in modo omogeneo, come se esistesse un unico $\varepsilon>0$ valido per \emph{tutti} i punti del dominio.\\
\\
Inoltre, il limite di una successione di funzioni Riemann--integrabili può non essere Riemann--integrabile, anche quando le funzioni $f_n$ sono monotone e non negative.

\begin{esempio}
Consideriamo le funzioni $f_n : [0,1] \to \mathbb{R}$ definite da
\[
f_n(x) = 
\begin{cases}
1, & \text{se } x \in \left[0, \tfrac{1}{n}\right],\\[4pt]
0, & \text{altrimenti.}
\end{cases}
\]
Allora $f_n(x) \uparrow f(x)$ con
\[
f(x) = 
\begin{cases}
1, & \text{se } x = 0,\\[4pt]
0, & \text{se } x > 0.
\end{cases}
\]

Ogni $f_n$ è Riemann-integrabile e si ha
\[
\int_0^1 f_n(x)\,dx = \frac{1}{n} \to 0.
\]
Tuttavia, la funzione limite $f$ è nulla quasi ovunque, ma \emph{non è Riemann-integrabile}, poiché presenta una discontinuità in un insieme di misura nulla che però non permette di definire il limite dell’integrale come quello delle funzioni $f_n$ nel senso di Riemann.

Il teorema di Beppo Levi richiederebbe che
\[
\int_0^1 f_n(x)\,dx \uparrow \int_0^1 f(x)\,dx,
\]
ma qui l’integrale a destra non è nemmeno definito in senso di Riemann, mentre in senso di Lebesgue si ha $\int_0^1 f = 0$ e dunque il teorema risulta valido solo per l’integrale di Lebesgue.
\end{esempio}
\bigskip
\noindent Ad ora siamo riusciti a definire l'integrale di Lebesgue per funzioni positive. Per estendere la definizione a qualsiasi funzione, ci basta fare un semplice trick matematico:
\begin{definizione}
    Data $f:E\to\mathbb{R}$ definiamo:
    \begin{itemize}
        \item $f_+:=f \vee 0 \qquad \text{parte positiva di }f $
        \item $f_-:=-(f \vee 0) = -f\wedge 0 \qquad \text{parte negativa di }f $
    \end{itemize}
\end{definizione}

Grazie a questa definizione notiamo che sia la parte positiva che la parte negativa sono funzioni entrambe positive, ossia che $f_-, f_+:E\to[0, +\infty]$. Inoltre possiamo riscrivere $f$ attraverso queste due nuove funzioni:
\[f = f_+-f_-\qquad \text{e anche che } \qquad |f| = f_++f_-\]

Siccome $f_-, f_+$ sono entrambe funzioni positive, allora esistono gli integrali $\int_Ef_+\;d\mu, \;\;\int_Ef_-\;d\mu$ definiti secondo la Definizione \ref{def:6.5}; potrebbero però valere $+\infty$.

\section{Spazio $\mathcal{L}^1$}
Siccome possiamo integrare sia la parte positiva che la parte negativa, ma questi potrebbero divergere a $\infty$, abbiamo la necessità di distinguere quando una funzione \textit{ammette} integrale, e quando invece la funzione è \textit{integrabile}.

\begin{definizione}[Ammissione dell'integrale]
    Data $f:E\to\mathbb{R}$, si dice che $f$ ammette integrale se 
    \[\int_Ef_+\:d\mu<+\infty\qquad \text{oppure}\qquad \int_Ef_-\:d\mu<+\infty\]
    In tal caso: 
    \[\int_Ef\:d\mu = \int_Ef_+\:d\mu-\int_Ef_-\:d\mu = 
    \begin{cases}
        +\infty\\
        -\infty<c<+\infty\\
        -\infty
    \end{cases}\]
\end{definizione}

\begin{definizione}[Funzione Integrabile]
\label{def:6.9}
    Data $f:E\to\mathbb{R}$ si dice integrabile se
    \[\int_Ef_+\:d\mu<+\infty\qquad \text{e}\qquad \int_Ef_-\:d\mu<+\infty\]
    In tal caso: 
    \[\boxed{\int_Ef\:d\mu = \int_Ef_+\:d\mu-\int_Ef_-\:d\mu} \in (-\infty, +\infty) \]
\end{definizione}

\begin{definizione}
    Definiamo lo spazio $\mathcal{L}^1$ come: 
    \[\mathcal{L}^1(E, \mathcal{E, \mu}) = \set{f:E\to\mathbb{R}\mid \;f\;\text{ è integrabile nel senso della Definizione \ref{def:6.9}}}\]
\end{definizione}

\noindent Date queste tre definizioni fondamentali, possiamo alcune osservazioni:
\begin{itemize}
    \item $\mathcal{L}^1(E, \mathcal{E}, \mu)$ è uno spazio lineare; questa proprietà discende direttamente dalla linearità dell'integrale di Lebesgue:
    \[a, b, \in \mathbb{R};\;f, g\in \mathcal{L}^1(E, \mathcal{E}, \mu) \quad \Rightarrow \quad af+bg\in \mathcal{L}^1(E, \mathcal{E}, \mu)\]
    \[\text{perchè: }\quad \int_E(af+bg)\;d\mu = a\int_Ef\;d\mu + b\int_Eg\;d\mu\]
    \item C'è coerenza fra tutte e tre le definizioni di integrale: \ref {def:intleb} $\equiv$ \ref{def:6.5} $\equiv$ \ref{def:6.9}
    \item Se due funzioni sono uguali quasi ovunque, allora anche queste avranno stesso integrale:
    \[f, g\in\mathcal{L}^1; \;\mu(f\neq g) = 0 \quad \Rightarrow \quad \int_Ef\:d\mu = \int_Eg\:d\mu\]
\end{itemize}

\begin{teorema}[Lebesgue o Convergenza Dominata]
    Se sono verificate le condizioni:
    \begin{enumerate}
        \item $f_n(x) \underset{n\to+\infty}{\longrightarrow} f(x)\quad \forall x \in E$
        \item $\exists \;g\in \mathcal{L}^1\;\text{ tale che }\;|f_n| \leq g\quad\forall n\geq0$
    \end{enumerate}
    Allora si ha che:

    \[f, f_n \in \mathcal{L}^1 \qquad \text{e}\qquad \underset{n\to+\infty}{\lim} \int_E f_n\;d\mu = \int_Ef\;d\mu\]
\end{teorema}

In altre parole, la condizione di dominanza da parte di $g$ permette di scambiare limite ed integrale in modo sicuro.\\
\\
Si presti particolare attenzione al fatto che si possono utilizzare altre scritture per la formula vista nel Teorema di Lebsegue:

\[\underset{n\to+\infty}{\lim} \int_E f_n\;d\mu = \int_Ef\;d\mu\quad\iff \quad \lim_{n\to+\infty}\left|\int_Ef_n\;d\mu - \int_Ef\;d\mu\right| = 0 \quad\iff \quad \int_E \left|f_n-f\;d\mu\right| \underset{n\to+\infty}{\longrightarrow}0\]

\bigskip
\noindent Fino ad ora abbiamo visto sempre l'integrale calcolato su tutto il dominio $E$, ma potremmo anche calcolarlo su sottoinsiemi di $E$ misurabili attraverso la seguente formuala:

\begin{definizione}
\label{def:6.12}
    $\forall A\in\mathcal{E}$ si ha che:
    \[\int_A f\;d\mu = \int_E f \cdot \mathbb{I}_{A}\;d\mu\qquad \forall f\in\mathcal{L}^1\]
\end{definizione}

\bigskip
\section{Caso di Interesse con $E = \mathbb{R}, \mathbb{R}^n$ e Confronto con Integrale di Reimann}

Tutta la teoria sviluppata fino ad ora è estremamente generale, tuttavia nelle applicazioni pratiche (o quantomeno per questo corso) si tenderà ad avere un caso specifico con:
\begin{itemize}
    \item $E = \mathbb{R}, \mathbb{R}^n$
    \item $\mu = m, $ misura di Lebesgue su $(\mathbb{R}, \mathcal{B}(\mathbb{R}))$ oppure $(\mathbb{R}^2, \mathcal{B}(\mathbb{R}^2))$.
\end{itemize}

\begin{esempio}
    Prendiamo la Definizione \ref{def:6.12} ed applichiamola al caso di interesse.\\
    Sia \(f(x) = c \, \mathbb{I}_{[a,b]}(x), \qquad c \geq 0.\)
    
\begin{center}
\begin{tikzpicture}[x=0.8cm,y=0.8cm,>=stealth]
    % Assi
    \draw[->,thick] (-1,0) -- (5,0);
    \draw[->,thick] (0,0) -- (0,2);
    % Segmento costante
    \draw[red,very thick] (1,1.2) -- (4,1.2);
    \draw[red,very thick] (-0.96,0) -- (1,0);
    \draw[red,very thick] (4,0) -- (4.96,0);
    % Linee tratteggiate
    \draw[dashed] (1,0) -- (1,1.2) node[below=1cm]{$a$};
    \draw[dashed] (4,0) -- (4,1.2) node[below=1cm]{$b$};
    % Etichetta
    \node[left] at (0,1.2) {$c$};
\end{tikzpicture}
\end{center}
La funzione $f$ è elementare e perciò integrabile:
\[
\int_{[a,b]} f \, d\mu 
= \int_{\mathbb{R}} c\, \mathbb{I}_{[a,b)}(x)\, dm
= 0 \cdot m(\mathbb{R}\setminus [a,b]) 
+ c \cdot m([a,b))
= c(b-a)
= c \int_a^b dx.
\]
\(\Rightarrow \quad f \text{ è Riemann--integrabile.}\)
\end{esempio}

\bigskip
\noindent Dall'esempio notiamo che può esistere una relazione tra l'integrale di Lebesgue e l'integrale di Reimann.

\begin{proposizione}
\label{prop:6.13}
    Se $\;h:[a, b]\to \mathbb{R}\;$ limitata e Riemann--integrabile \(\quad \Rightarrow\quad h \; \) è Lebesgue--integrabile rispetto a $m$ su $\mathcal{B}$ e vale che i due integrali coincidono:
    \[\int_a^bh(x)\;dx = \int_{[a, b]} h \; dm\]
\end{proposizione}

\bigskip
\noindent Osserviamo che possiamo scrivere:
\[|h| = h_+ + h_-\qquad \Rightarrow \qquad \int_\mathbb{R}|h|\;dm =\int_\mathbb{R}h_+\;dm + \int_\mathbb{R}h_-\;dm \]

Di conseguenza, se $h$ è integrabile, allora anche $|h|$ è integrabile. Questo è valido solo per l'integrazione secondo Lebesgue, in quanto siamo sicuri che $h_+$ e $h_-$ sono integrabili se lo è $h$. Per quanto riguarda l'integrazione di Reimann non possiamo decomporre $h$ nella differenza tra la sua parte positiva e negativa ed integrare; questo perché se una funzione $h$ è integrabile, allora non è detto che $|h|$ sia integrabile\footnote{Si prenda il classico esempio di $\frac{\sin{x}}{x}$}.\\
Se invece ipotizziamo che $|h|$ sia integrabile:

\begin{align*}
    \left| \int_\mathbb{R} h\;dm\right| & = \left|\int_\mathbb{R}h_+\;dm - \int_\mathbb{R}h_-\;dm  \right| \leq \left|\int_\mathbb{R}h_+\;dm \right|+ \left|\int_\mathbb{R}h_-\;dm  \right| \\
    & = \int_\mathbb{R}h_+\;dm + \int_\mathbb{R}h_-\;dm\\
    & = \int_\mathbb{R} |h|\;dm < +\infty
\end{align*}

Unendo queste due condizioni che abbiamo appena ricavato, \textbf{scopriamo una proprietà fondamentale dell'integrale di Lebesgue, che non vale per quello di Riemann}.
\[\boxed{h\in\mathcal{L}^1(\mathbb{R}, \mathcal{B}, m) \quad \iff \quad |h| \in \mathcal{L}^1(\mathbb{R}, \mathcal{B}, m)}\]
Possiamo quindi ridefinire lo spazio $\mathcal{L}^1$ come: $\quad \boxed{\mathcal{L}^1(E, \mathcal{E}, \mu) = \set{f:E\to\mathbb{R}\big{|}\; |f|\; \text{ è integrabile}}}$

\newpage
\noindent \textit{È quindi corretto vedere l'integrale di Lebesgue come una ``generalizzazione'' dell'integrale di Reimann?}

\medskip
\noindent
\textbf{In un certo senso sì, per gli integrali propri}. Infatti, ogni funzione Riemann--integrabile su $[a,b]$ è anche Lebesgue--integrabile, e i due integrali coincidono, come abbiamo visto con la Proposizione \ref{prop:6.13}.

\medskip
L'integrale di Lebesgue è però definito per un insieme molto più ampio di funzioni, poiché si basa sulla misura degli insiemi di punti in cui la funzione assume determinati valori, piuttosto che sulla suddivisione del dominio in intervalli come nel caso di Riemann. In questo modo, l'integrale di Lebesgue risulta adatto a trattare funzioni non limitate, con infiniti punti di discontinuità, o definite su spazi di misura generali, pur mantenendo le proprietà fondamentali di linearità e monotonia dell'integrale classico.

\begin{esempio}
    Prendiamo la funzione di Dirichlet $\mathbb{I}_{\mathbb{Q}}(x)$, che è ben nota per non essere Riemann--integrabile, tuttavia:
    \[\int_\mathbb{R} \mathbb{I}_{\mathbb{Q}}(x) \; dm = m(\mathbb{Q}) = m\left(\bigcup_n \{q_n\}\right) = \sum_{n\geq1}m(\{q_n\}) = 0\]
\end{esempio}

\bigskip
\noindent \textbf{Nel caso degli integrali impropri}, l'integrale di Lebesgue \textbf{rappresenta un'estensione} rigorosa dell'integrale di Riemann \textbf{solo nei casi di convergenza assoluta}. Infatti, se un integrale improprio di Riemann converge assolutamente, allora esso coincide con l'integrale di Lebesgue.

\medskip
\noindent Esistono casi dove però Lebesgue non ammette convergenza e Riemann si (in maniera impropria), per esempio in tutte quelle funzioni che hanno oscillazioni e ``compensazione di segno'': se $\int |f| = +\infty$, l'integrale non è definito, anche se quello improprio di Riemann converge. 

\begin{esempio}
    Un esempio classico è $f(x) = \tfrac{\sin x}{x}$ su $[0,+\infty)$, che è integrabile in senso di Riemann improprio, ma non in senso di Lebesgue.
    \label{leb_ndim}
\end{esempio}

\bigskip
\noindent Possiamo quindi concludere questa trattazione generalizzando e dicendo che:
\begin{itemize}
    \item Nel caso in cui $E = \mathbb{R}, \mathcal{E} = \mathcal{B}(\mathbb{R}), \mu = m$, allora se $f\geq0$ e Riemann--integrabile su illimitati $\Rightarrow\; f$ è Lebesgue integrabile;

    \item Nel caso in cui $E = \mathbb{R}^n, \mathcal{E} = \mathcal{B}(\mathbb{R}^n), \mu = m_n$, definiamo la misura di Lebesgue n--dimensionale come
    \[m_n = ([a_1, b_1)\times\cdots\times[a_n, b_n)) = \prod_{k=1}^n(b_k-a_k)\]
    e se $f$ è Riemann--integrabile e limitata su $\mathbb{R}^n$, allora $f$ è Lebesgue--integrabile, e i due integrali coincidono:
    \[\int_{([a_1, b_1)\times\cdots\times[a_n, b_n)}f(x_1, ..., x_n)\;dm_n = \int_{([a_1, b_1)\times\cdots\times[a_n, b_n)}f(x_1, ..., x_n)\;dx_1...dx_n\]
\end{itemize}

